% --- Criação de Variáveis Customizadas ---
\providecommand{\imprimiravaliador}{}
\newcommand{\avaliador}[1]{\renewcommand{\imprimiravaliador}{#1}}

\autor{MATHEUS COSTA SOARES DE AZEVEDO}
\titulo{DESENVOLVIMENTO DE UMA SOLUÇÃO DE \textit{SOFTWARE} EM LINGUAGEM C PARA SUPORTE À IMPLEMENTAÇÃO DE \textit{ASSET ADMINISTRATION SHELL} EMBARCADOS APLICADOS À INDUSTRIA 4.0}
\orientador{Rubens de Andrade Fernandes, Dr.}
\avaliador{Bruno da Gama Monteiro, Me.}

\local{Manaus}
\data{2026}


\preambulo{Projeto de pesquisa desenvolvido durante a disciplina de Trabalho de Conclusão de Curso I e apresentado à banca avaliadora do curso de Engenharia Eletrônica da Escola Superior de Tecnologia da Universidade do Estado do Amazonas, como pré-requisito para obtenção do título de Engenheiro em Eletrônica.}

% --- Configurações de Formatação de Parágrafo ABNT ---
\setlength{\parindent}{1.5cm} % Recuo da primeira linha do parágrafo
\setlength{\parskip}{0.2cm}   % Espaço entre um parágrafo e outro
\OnehalfSpacing               % Espaçamento de 1.5 entre linhas

% --- Estilo de Página (Cabeçalhos) ---
% Força o estilo 'abntheadings' a ser 'simple', o que remove a linha e o nome do capítulo
% do topo da página, deixando apenas o número da página no canto superior direito.
\aliaspagestyle{abntheadings}{simple}

% --- Configurações de Links do PDF ---
% Isso remove as caixas vermelhas em volta dos itens do sumário e deixa tudo em texto preto normal
\hypersetup{
    colorlinks=true,
    linkcolor=black,
    citecolor=black,
    filecolor=black,
    urlcolor=black
}

% --- Traduções e Textos ---
\addto\captionsbrazil{
  \renewcommand{\contentsname}{SUMÁRIO}
  \renewcommand{\bibname}{REFERÊNCIAS}
  \renewcommand{\listfigurename}{LISTA DE ILUSTRAÇÕES}
}
